\chapter{Conception technique}
\markboth{Chapitre 3}{Conception technique}
\setcounter{chapter}{3}
\setcounter{section}{0}
\minitoc
\newpage

%% ═══════════════════════════════════════════════════════════════════════════════
\section{Introduction}

Ce chapitre détaille la conception technique de chaque composant de la plateforme~:
l'architecture des collecteurs, le pipeline d'ingestion avec sa formule de scoring,
le module NLP, le générateur de messages LLM, le schéma de base de données, et
les workflows d'automatisation n8n.

%% ═══════════════════════════════════════════════════════════════════════════════
\section{Architecture des collecteurs}

\subsection{Collecteur LinkedIn Jobs — Signal d'externalisation}

\textbf{Fichier~:} \texttt{collectors/linkedin\_hiring.py} \\
\textbf{Signal produit~:} \texttt{outsourcing\_signal}

\textbf{Logique métier~:} une entreprise non-IT qui recrute activement un
développeur COBOL possède un projet COBOL actif avec budget. Solvinya peut
contacter le DSI pour proposer une équipe externalisée immédiatement disponible,
avant même que le recrutement aboutisse.

\begin{figure}[H]
  \centering
  \begin{tikzpicture}[
    node distance=0.5cm,
    proc/.style={draw=darkGreen!70, fill=colColor, rounded corners=4pt,
                 minimum width=8cm, minimum height=0.8cm,
                 font=\small, align=center, inner sep=6pt},
    ann/.style={font=\scriptsize\itshape\color{espritGray},
                text width=4cm, align=left},
    arr/.style={-{Stealth[length=5pt]}, darkGreen!70, thick}
  ]
    \node[proc] (loop1) {BOUCLE externe : technologie (COBOL, SAP, Java, \ldots)};
    \node[proc, below=0.4cm of loop1] (loop2) {BOUCLE interne : pays (FR, BE, LU, CH, TN)};
    \node[proc, below=0.4cm of loop2] (search) {\texttt{search\_jobs(keywords, location, limit=30)}};
    \node[ann, right=0.5cm of search] {Retourne stubs\\(entityUrn seulement)};
    \node[proc, below=0.4cm of search] (detail) {\texttt{get\_job(job\_id)} → détail complet};
    \node[proc, below=0.4cm of detail] (extract) {\texttt{\_extract\_company\_info()} → universalName, name};
    \node[proc, below=0.4cm of extract] (prospect) {\texttt{HiringProspect(company, tech, country, boost)}};
    \node[proc, below=0.4cm of prospect] (ingest) {\texttt{ingest\_opportunity(type="outsourcing\_signal")}};
    \node[proc, below=0.4cm of ingest, fill=lightAmber!40, draw=espritRed!40]
      (sleep) {\texttt{sleep(gauss(4.0, 1.5))} — anti-détection};

    \foreach \a/\b in {loop1/loop2, loop2/search, search/detail, detail/extract,
                       extract/prospect, prospect/ingest, ingest/sleep}
      \draw[arr] (\a) -- (\b);
  \end{tikzpicture}
  \caption{Flux interne du collecteur LinkedIn Jobs}
  \label{fig:flow_linkedin_hiring}
\end{figure}

\subsection{Collecteur LinkedIn B2B — Signaux directs}

\textbf{Fichier~:} \texttt{collectors/linkedin\_b2b.py} \\
\textbf{Signaux produits~:} \texttt{b2b\_rfp}, \texttt{b2b\_tender},
  \texttt{b2b\_subcontracting}, \texttt{b2b\_partnership}

\textbf{Contrainte technique~:} l'API LinkedIn non-officielle ne retourne aucun
résultat avec le filtre \texttt{CONTENT} sur les posts. La stratégie adoptée est
donc indirecte~: rechercher les \emph{entreprises acheteuses} par secteur, puis
scanner leurs publications récentes.

\begin{figure}[H]
  \centering
  \begin{tikzpicture}[
    node distance=0.45cm,
    proc/.style={draw=darkBlue!70, fill=pipColor, rounded corners=4pt,
                 minimum width=9cm, minimum height=0.8cm,
                 font=\small, align=center, inner sep=6pt},
    bad/.style={draw=espritRed!60, fill=warnColor!50, rounded corners=4pt,
                minimum width=9cm, minimum height=0.7cm,
                font=\small\itshape, align=center},
    good/.style={draw=darkGreen!70, fill=colColor, rounded corners=4pt,
                 minimum width=9cm, minimum height=0.8cm,
                 font=\small, align=center, inner sep=6pt},
    arr/.style={-{Stealth[length=5pt]}, darkBlue!70, thick},
    garr/.style={-{Stealth[length=5pt]}, darkGreen!70, thick}
  ]
    \node[bad] (bad1)
      {\texttt{search\_content("appel d'offres Java")} → \textbf{0 résultats} $\times$};

    \node[proc, below=0.6cm of bad1] (q1)
      {\texttt{search\_companies("banque crédit financement France", limit=8)}};
    \node[proc, below=0.4cm of q1] (q2)
      {Pour chaque entreprise → \texttt{get\_company\_updates\_by\_id(urn\_id, 15)}};
    \node[proc, below=0.4cm of q2] (q3)
      {\texttt{\_extract\_post\_text(update)} → texte brut};
    \node[proc, below=0.4cm of q3] (filter)
      {Filtre \texttt{\_ESN\_INDICATORS} → exclure les ESN auto-promotionnels};
    \node[good, below=0.4cm of filter] (score)
      {\texttt{\_score\_post(text)} → \texttt{(opp\_type, technology, boost)}};
    \node[good, below=0.4cm of score] (ingest)
      {\texttt{ingest\_raw\_dict(sig, source\_platform="linkedin\_b2b")}};

    \draw[arr, dashed, espritRed!60] (bad1) -- node[right, font=\scriptsize]
      {Limitation API} (q1);
    \foreach \a/\b in {q1/q2, q2/q3, q3/filter, filter/score, score/ingest}
      \draw[arr] (\a) -- (\b);
  \end{tikzpicture}
  \caption{Stratégie et flux du collecteur LinkedIn B2B}
  \label{fig:flow_linkedin_b2b}
\end{figure}

Les patterns de détection regex sont détaillés dans le tableau~\vref{tab:b2b_patterns}.

\begin{table}[H]
  \centering
  \begin{tabular}{L{5.5cm} L{3.5cm} C{1.5cm}}
    \toprule
    \rowcolor{lightBlue!50}
    \textbf{Expression régulière} & \textbf{Type signal} & \textbf{Boost} \\
    \midrule
    \texttt{rfp|request for proposal|ddp} & \texttt{b2b\_rfp} & $+22$ \\
    \texttt{appel d.offres|marché public|consultation} & \texttt{b2b\_tender} & $+18$ \\
    \texttt{recherche prestataire|sous-traitan|outsourc} & \texttt{b2b\_subcontracting} & $+16$ \\
    \texttt{partenariat stratégique|co-traitan|consortium} & \texttt{b2b\_partnership} & $+14$ \\
    \texttt{sélectionné|retenu|choisi + prestataire} & \texttt{b2b\_tender} & $+12$ \\
    \bottomrule
  \end{tabular}
  \caption{Patterns regex de détection B2B et leurs scores de boost}
  \label{tab:b2b_patterns}
\end{table}

\subsection{Collecteur BOAMP}

\textbf{Fichier~:} \texttt{collectors/boamp.py} |
\textbf{API~:} \texttt{boamp-datadila.opendatasoft.com/api/explore/v2.1}

Le BOAMP (Bulletin Officiel des Annonces de Marchés Publics) expose une API REST
open data gratuite. Pour chacun des 15 mots-clés IT, une requête GET filtre les
avis dont l'objet contient le terme, triés par date de parution décroissante.

\begin{lstlisting}[language=Python3, caption={Requête BOAMP (extrait)}, label=lst:boamp]
GET /records
  ?where=objet like "COBOL"
  &limit=20
  &order_by=dateparution desc
  &select=objet,nomacheteur,url_avis,dateparution,nature_libelle
\end{lstlisting}

\subsection{Collecteur TED EU}

\textbf{Fichier~:} \texttt{collectors/ted.py} |
\textbf{API~:} \texttt{api.ted.europa.eu/v3/notices/search} (POST uniquement)

TED (Tenders Electronic Daily) couvre les marchés publics de l'UE. La syntaxe
\emph{expert query} permet de filtrer simultanément par titre et par pays.

\begin{lstlisting}[language=Python3, caption={Requête TED EU (expert query)}, label=lst:ted]
POST https://api.ted.europa.eu/v3/notices/search
{
  "query": "notice-title=\"SAP\" AND buyer-country=LUX",
  "fields": ["notice-title", "buyer-name",
             "publication-date", "notice-url", "cpv-codes"],
  "limit": 15
}
\end{lstlisting}

Points critiques de conception~:
\begin{itemize}
  \item Chaque requête couvre \textbf{un seul pays} (boucle Python sur FRA, BEL, LUX, CHE) ;
  \item Le nom de l'acheteur est un dictionnaire multilingue~: \texttt{buyer-name.get("fra") or .get("eng")} ;
  \item Les codes pays utilisent la norme ISO alpha-3 : \texttt{FRA}, \texttt{BEL}, \texttt{LUX}, \texttt{CHE}.
\end{itemize}

%% ═══════════════════════════════════════════════════════════════════════════════
\section{Pipeline d'ingestion}

\subsection{Vue d'ensemble — 6 étapes, $\approx$50\,ms/prospect}

\begin{figure}[H]
  \centering
  \begin{tikzpicture}[
    pipstep/.style={draw=darkBlue!60, fill=pipColor, rounded corners=5pt,
                 minimum width=11.5cm, minimum height=0.9cm,
                 font=\small\bfseries, align=center, inner sep=8pt},
    num/.style={draw=espritBlue, fill=espritBlue, circle,
                minimum size=0.6cm, text=white, font=\small\bfseries},
    arr/.style={-{Stealth[length=5pt]}, darkBlue!60, thick}
  ]
    \foreach \n/\txt/\sub [count=\i from 0] in {
      1/{Normalisation des champs}/{country\_iso → ISO 2 lettres · tech → taxonomie 50+ · sector → fixe},
      2/{UPSERT entreprise}/{linkedin\_url comme clé unique → retourne company\_id (UUID)},
      3/{Déduplication 30 jours}/{même entreprise + même tech + même type dans fenêtre 30j → SKIP},
      4/{Calcul du score 0--100}/{\texttt{compute\_rule\_score(country, tech, sector, boost)}},
      5/{Génération du titre contextuel}/{\textit{[Appel d'offres] SAP — Gouvernement LU}},
      6/{INSERT opportunité}/{status='new' · priority\_score · source\_platform · technologies[]}
    }{
      \node[pipstep] (s\i) at (0, {-\i*1.2}) {Étape \n~: \txt \\ {\scriptsize\normalfont\color{espritGray}\sub}};
      \ifnum \i>0
        \draw[arr] ($(s\the\numexpr\i-1\relax.south)$) -- ($(s\i.north)$);
      \fi
    }
  \end{tikzpicture}
  \caption{Les 6 étapes du pipeline \texttt{ingest.py}}
  \label{fig:pipeline}
\end{figure}

\subsection{Formule de scoring — 4 dimensions pondérées}

\begin{infobox}[Formule de scoring]
\[
  \text{score} = \underbrace{S_{\text{pays}} \times 0.30}_{\text{30\%}}
               + \underbrace{S_{\text{tech}} \times 0.25}_{\text{25\%}}
               + \underbrace{S_{\text{secteur}} \times 0.25}_{\text{25\%}}
               + \underbrace{\frac{\text{boost}}{25} \times 100 \times 0.20}_{\text{20\%}}
\]
où chaque $S_i \in [0, 100]$. Le résultat est plafonné à $100$.
\end{infobox}

\begin{figure}[H]
  \centering
  \begin{tikzpicture}
    \begin{axis}[
      xbar stacked,
      width=13cm, height=4.5cm,
      bar width=18pt,
      xmin=0, xmax=100,
      xlabel={Score / 100},
      symbolic y coords={React-FR, Python-CH, SAP-BE, COBOL-LU},
      ytick=data,
      y tick label style={font=\small},
      legend style={at={(0.5,-0.35)}, anchor=north, font=\small,
                    legend columns=4},
      enlarge y limits=0.3,
      nodes near coords,
      nodes near coords align={horizontal},
      every node near coord/.style={font=\scriptsize\color{white}\bfseries},
    ]
      \addplot+[fill=espritBlue!70] coordinates
        {(27,{COBOL-LU}) (25.5,{SAP-BE})
         (28.5,{Python-CH}) (27,{React-FR})};
      \addplot+[fill=darkGreen!70] coordinates
        {(25,{COBOL-LU}) (21.25,{SAP-BE})
         (17.5,{Python-CH}) (13.75,{React-FR})};
      \addplot+[fill=espritRed!70] coordinates
        {(25,{COBOL-LU}) (17.5,{SAP-BE})
         (18.75,{Python-CH}) (16.25,{React-FR})};
      \addplot+[fill=lightAmber!80, draw=espritRed!40] coordinates
        {(20,{COBOL-LU}) (9.6,{SAP-BE})
         (12,{Python-CH}) (4,{React-FR})};
      \legend{Pays (30\%), Tech (25\%), Secteur (25\%), Boost (20\%)}
    \end{axis}
  \end{tikzpicture}
  \caption{Décomposition du score pour 4 exemples types}
  \label{fig:score_decomp}
\end{figure}

\begin{table}[H]
  \centering
  \begin{tabular}{C{1.5cm} C{2cm} C{2cm} C{2.5cm} C{1.8cm} C{1.5cm}}
    \toprule
    \rowcolor{lightBlue!50}
    \textbf{Pays} & \textbf{$S_{\text{pays}}$} &
    \textbf{Technologie} & \textbf{$S_{\text{tech}}$} &
    \textbf{Secteur} & \textbf{$S_{\text{sect}}$} \\
    \midrule
    Luxembourg & 100 & COBOL & 100 & Finance/Banque & 100 \\
    Suisse & 95 & Mainframe & 100 & Finance/Assur. & 95 \\
    France & 90 & SAP & 85 & Finance & 90 \\
    Belgique & 85 & Java & 80 & IA & 75 \\
    Tunisie & 65 & ML & 78 & Industrie/RH & 70 \\
    & & Python/DevOps & 70 & Tech/IA & 65 \\
    & & Salesforce & 65 & Retail & 50 \\
    & & React & 55 & Secteur public & 50 \\
    \bottomrule
  \end{tabular}
  \caption{Table des scores par dimension}
  \label{tab:score_dims}
\end{table}

%% ═══════════════════════════════════════════════════════════════════════════════
\section{Module NLP}

\subsection{Architecture des deux classifieurs}

Le module NLP (\texttt{nlp/classifier.py}) contient deux pipelines
scikit-learn indépendants, illustrés en figure~\vref{fig:nlp_arch}.

\begin{figure}[H]
  \centering
  \begin{tikzpicture}[
    box/.style={draw=#1!70, fill=#1!15, rounded corners=4pt,
                minimum width=4.8cm, minimum height=0.8cm,
                font=\small, align=center, inner sep=6pt},
    arr/.style={-{Stealth[length=5pt]}, #1!70, thick},
    node distance=0.5cm
  ]
    %% ─ Classifieur Secteur ────────────────────────────────────────────────────
    \node[box=espritBlue, minimum width=5cm, minimum height=0.6cm,
          font=\small\bfseries] (in1) at (-3, 6)   {Texte brut FR};
    \node[box=darkBlue] (tf1a) at (-3, 5)  {TF-IDF word 1--3g\\$(n\_features=50000)$};
    \node[box=darkBlue] (tf1b) at (-3, 3.9){TF-IDF char 3--5g\\$(n\_features=30000)$};
    \node[box=darkBlue] (fu1)  at (-3, 2.8){FeatureUnion → concaténation};
    \node[box=nlpColor, draw=espritRed!60] (svc1) at (-3, 1.7)
      {LinearSVC + CalibratedClassifierCV};
    \node[box=darkGreen] (out1) at (-3, 0.6)
      {\texttt{sector\_label} + confiance};

    \node[draw=espritBlue!40, dashed, rounded corners=6pt,
          fit=(in1)(tf1a)(tf1b)(fu1)(svc1)(out1),
          label={[font=\small\bfseries\color{espritBlue}]above:Classifieur Secteur (10 classes)}]
      {};

    %% ─ Classifieur Priorité ──────────────────────────────────────────────────
    \node[box=espritRed, minimum width=5cm, minimum height=0.6cm,
          font=\small\bfseries] (in2) at (3.5, 6)   {Texte brut FR};
    \node[box=darkBlue] (tf2a) at (3.5, 5)  {TF-IDF word 1--3g};
    \node[box=darkBlue] (tf2b) at (3.5, 3.9){TF-IDF char 3--5g};
    \node[box=darkBlue] (fu2)  at (3.5, 2.8){FeatureUnion → concaténation};
    \node[box=nlpColor, draw=espritRed!60] (svc2) at (3.5, 1.7)
      {LinearSVC + CalibratedClassifierCV};
    \node[box=darkGreen] (out2) at (3.5, 0.6)
      {\texttt{priority\_label} HIGH/MED/LOW};

    \node[draw=espritRed!40, dashed, rounded corners=6pt,
          fit=(in2)(tf2a)(tf2b)(fu2)(svc2)(out2),
          label={[font=\small\bfseries\color{espritRed}]above:Classifieur Priorité (3 classes)}]
      {};

    %% Flèches
    \foreach \a/\b in {in1/tf1a, tf1a/tf1b, tf1b/fu1, fu1/svc1, svc1/out1}
      \draw[arr=darkBlue] (\a) -- (\b);
    \foreach \a/\b in {in2/tf2a, tf2a/tf2b, tf2b/fu2, fu2/svc2, svc2/out2}
      \draw[arr=darkBlue] (\a) -- (\b);
  \end{tikzpicture}
  \caption{Architecture des deux pipelines NLP (TF-IDF + LinearSVC)}
  \label{fig:nlp_arch}
\end{figure}

\subsection{Comparaison des modèles}

\begin{figure}[H]
  \centering
  \begin{tikzpicture}
    \begin{axis}[
      ybar,
      width=11cm, height=5.5cm,
      bar width=28pt,
      ymin=60, ymax=82,
      ylabel={Accuracy CV-5 (\%)},
      symbolic x coords={TF-IDF (word+char n-grams), Logistic Regression, sentence-transformers},
      xtick=data,
      x tick label style={font=\small, align=center, text width=3.5cm},
      nodes near coords,
      nodes near coords style={font=\small\bfseries},
      enlarge x limits=0.3,
    ]
      \addplot+[fill=darkGreen!70, draw=darkGreen] coordinates
        {(TF-IDF (word+char n-grams), 75.0)};
      \addplot+[fill=espritGray!50, draw=espritGray] coordinates
        {(Logistic Regression, 71.0)};
      \addplot+[fill=espritBlue!50, draw=espritBlue] coordinates
        {(sentence-transformers, 69.4)};
    \end{axis}
    %% Ligne seuil
    \draw[espritRed, dashed, thick] (1.2, 1.65) -- (9.3, 1.65)
      node[right, font=\small\bfseries\color{espritRed}] {Seuil 75\%};
  \end{tikzpicture}
  \caption{Comparaison des approches NLP — corpus de 604 exemples IT}
  \label{fig:nlp_comparison}
\end{figure}

\subsection{Distribution des données d'entraînement}

\begin{figure}[H]
  \centering
  \begin{tikzpicture}
    \begin{axis}[
      xbar,
      width=10cm, height=6cm,
      bar width=14pt,
      xmin=0, xmax=140,
      xlabel={Nombre d'exemples},
      symbolic y coords={Autres, Transport-Retail, Santé,
                         Secteur-public, Industrie-RH, Tech-IA, Finance-Banque},
      ytick=data,
      y tick label style={font=\small},
      nodes near coords,
      nodes near coords align={horizontal},
      enlarge y limits=0.12,
    ]
      \addplot+[fill=espritBlue!70] coordinates {
        (54,{Autres})
        (60,{Transport-Retail})
        (70,{Santé})
        (90,{Secteur-public})
        (100,{Industrie-RH})
        (110,{Tech-IA})
        (120,{Finance-Banque})
      };
    \end{axis}
  \end{tikzpicture}
  \caption{Répartition des 604 exemples d'entraînement NLP par secteur}
  \label{fig:nlp_data}
\end{figure}

%% ═══════════════════════════════════════════════════════════════════════════════
\section{Générateur de messages LLM}

\begin{figure}[H]
  \centering
  \begin{tikzpicture}[
    node distance=0.5cm,
    box/.style={draw=#1!60, fill=#1!15, rounded corners=4pt,
                minimum width=11cm, minimum height=0.8cm,
                font=\small, align=center, inner sep=6pt},
    arr/.style={-{Stealth[length=5pt]}, darkBlue!60, thick}
  ]
    \node[box=espritBlue] (api) {\texttt{POST /api/generate/message} · \{"opportunity\_id": "uuid"\}};
    \node[box=darkGreen, below=0.4cm of api] (read)
      {Lecture de l'opportunité en DB (company, tech, type, score, sector)};
    \node[box=darkBlue, below=0.4cm of read] (prompt)
      {Construction du prompt adapté au \texttt{opportunity\_type}};

    \node[draw=darkBlue!30, rounded corners=4pt, fill=lightBlue!20,
          below=0.4cm of prompt, minimum width=11cm, inner sep=8pt] (ptypes) {
      \begin{tabular}{ll}
        \texttt{outsourcing\_signal} & → \textit{"L'entreprise recrute un dev \{tech\}. Proposer l'externalisation."} \\
        \texttt{b2b\_tender}         & → \textit{"AO IT détecté. Proposer une réponse complète."} \\
        \texttt{b2b\_partnership}    & → \textit{"Cherche partenaire. Présenter Solvinya."} \\
        \texttt{b2b\_rfp}            & → \textit{"RFP publié. Proposer une solution complète."} \\
      \end{tabular}
    };

    \node[box=espritBlue, below=0.4cm of ptypes]
      (groq) {Tentative 1~: \textbf{Groq API} (\texttt{llama-3.1-8b-instant}) — latence 1--3s};
    \node[box=espritGray, below=0.3cm of groq]
      (ollama) {Fallback~: \textbf{Ollama local} (\texttt{llama3:8b}) — latence 5--15s};
    \node[box=darkGreen, below=0.4cm of ollama]
      (out) {Message LinkedIn 5 lignes · INSERT dans \texttt{leads}};

    \foreach \a/\b in {api/read, read/prompt, prompt/ptypes, ptypes/groq, groq/ollama, ollama/out}
      \draw[arr] (\a) -- (\b);
    %% Séparateur Groq/Ollama
    \draw[espritRed!50, dashed] (-5.6, -4.55) -- (5.5, -4.55)
      node[right, font=\scriptsize\color{espritRed}] {Si Groq indisponible};
  \end{tikzpicture}
  \caption{Pipeline de génération de messages via LLM}
  \label{fig:llm_pipeline}
\end{figure}

%% ═══════════════════════════════════════════════════════════════════════════════
\section{Schéma de base de données}

\begin{figure}[H]
  \centering
  \begin{tikzpicture}[
    tbl/.style={draw=espritBlue!60, rounded corners=4pt, fill=lightBlue!20},
    hdr/.style={fill=espritBlue!70, text=white, font=\small\bfseries},
    row/.style={font=\scriptsize, minimum height=0.5cm},
    pk/.style={font=\scriptsize\bfseries\color{espritRed}},
    fk/.style={font=\scriptsize\itshape\color{darkGreen}},
    rel/.style={-{Stealth[length=5pt]}, espritGray!70, thick}
  ]
    %% ── Table companies ───────────────────────────────────────────────────────
    \matrix (companies) [tbl, matrix of nodes,
      nodes in empty cells, column sep=0pt, row sep=0pt,
      nodes={draw=espritBlue!20, minimum width=3.5cm, inner sep=4pt, anchor=west}]
      at (-4, 3) {
        |[hdr, minimum width=3.5cm]| companies \\
        |[pk]| id : UUID PK \\
        |[row]| name : VARCHAR(255) \\
        |[row]| linkedin\_url : UNIQUE \\
        |[row]| country : CHAR(2) \\
        |[row]| sector : VARCHAR(80) \\
        |[row]| source : VARCHAR(50) \\
        |[row]| created\_at : TIMESTAMPTZ \\
      };

    %% ── Table opportunities ───────────────────────────────────────────────────
    \matrix (opps) [tbl, matrix of nodes,
      nodes in empty cells, column sep=0pt, row sep=0pt,
      nodes={draw=espritBlue!20, minimum width=4.5cm, inner sep=4pt, anchor=west}]
      at (1, 3) {
        |[hdr, minimum width=4.5cm]| opportunities \\
        |[pk]| id : UUID PK \\
        |[fk]| company\_id : UUID FK \\
        |[row]| title : VARCHAR(500) \\
        |[row]| opportunity\_type : VARCHAR(50) \\
        |[row]| technologies : TEXT{[}{]} \\
        |[row]| source\_platform : VARCHAR(80) \\
        |[row]| country : CHAR(2) \\
        |[row]| priority\_score : SMALLINT \\
        |[row]| status : VARCHAR(30) \\
        |[row]| created\_at : TIMESTAMPTZ \\
      };

    %% ── Table contacts ────────────────────────────────────────────────────────
    \matrix (contacts) [tbl, matrix of nodes,
      nodes in empty cells, column sep=0pt, row sep=0pt,
      nodes={draw=espritBlue!20, minimum width=3.5cm, inner sep=4pt, anchor=west}]
      at (-4, -2) {
        |[hdr, minimum width=3.5cm]| contacts \\
        |[pk]| id : UUID PK \\
        |[fk]| company\_id : UUID FK \\
        |[row]| full\_name : VARCHAR(200) \\
        |[row]| job\_title : VARCHAR(200) \\
        |[row]| email : VARCHAR(200) \\
        |[row]| linkedin\_url : VARCHAR(500) \\
        |[row]| is\_decision\_maker : BOOL \\
      };

    %% ── Table leads ───────────────────────────────────────────────────────────
    \matrix (leads) [tbl, matrix of nodes,
      nodes in empty cells, column sep=0pt, row sep=0pt,
      nodes={draw=espritBlue!20, minimum width=3.5cm, inner sep=4pt, anchor=west}]
      at (2, -2) {
        |[hdr, minimum width=3.5cm]| leads \\
        |[pk]| id : UUID PK \\
        |[fk]| opportunity\_id : UUID FK \\
        |[fk]| contact\_id : UUID FK \\
        |[row]| generated\_linkedin\_msg \\
        |[row]| generated\_email \\
        |[row]| status : VARCHAR(30) \\
        |[row]| sent\_at : TIMESTAMPTZ \\
      };

    %% Relations
    \draw[rel] (companies.east |- companies-2-1) -- ++(0.3,0)
      |- (opps.west |- opps-3-1);
    \draw[rel] (companies.east |- companies-2-1) -- ++(0.3,0)
      |- (contacts.west |- contacts-3-1);
    \draw[rel] (opps.south |- opps-2-1) -- ++(0,-0.3)
      |- (leads.east |- leads-3-1);
    \draw[rel] (contacts.east |- contacts-2-1) -- ++(0.3,0)
      |- (leads.west |- leads-4-1);
  \end{tikzpicture}
  \caption{Schéma entité-relation de la base de données}
  \label{fig:erd}
\end{figure}

%% ═══════════════════════════════════════════════════════════════════════════════
\section{Conception des workflows n8n}

\subsection{Workflow 04 — Déclencheur collecte (principal)}

Ce workflow est le plus important en production. Il déclenche la collecte complète
toutes les 12 heures via l'API, et permet également un déclenchement manuel par webhook.

\begin{figure}[H]
  \centering
  \begin{tikzpicture}[
    trig/.style={draw=darkGreen!70, fill=colColor, rounded corners=5pt,
                 minimum width=3cm, minimum height=1.1cm,
                 font=\small\bfseries, align=center},
    proc/.style={draw=darkBlue!70, fill=pipColor, rounded corners=5pt,
                 minimum width=3.5cm, minimum height=1.1cm,
                 font=\small, align=center},
    cond/.style={draw=espritRed!60, fill=warnColor!40, diamond,
                 minimum width=3cm, minimum height=1cm,
                 font=\small\bfseries, align=center, aspect=2},
    ok/.style={draw=darkGreen!70, fill=outColor, rounded corners=5pt,
               minimum width=3cm, minimum height=0.9cm,
               font=\small, align=center},
    err/.style={draw=espritRed!60, fill=warnColor!30, rounded corners=5pt,
                minimum width=3cm, minimum height=0.9cm,
                font=\small, align=center},
    arr/.style={-{Stealth[length=5pt]}, espritGray!70, thick}
  ]
    %% Déclencheurs
    \node[trig] (sched) at (-2, 4.5) {[Sch.] Schedule\\Toutes les 12h};
    \node[trig] (wh) at (-2, 2.5) {[WH] Webhook\\POST /webhook/collect};

    %% HTTP
    \node[proc] (http) at (3, 3.5) {[HTTP] HTTP Request\\POST /api/collect/all\\X-API-Key · timeout 10s};

    %% IF
    \node[cond] (cond) at (7, 3.5) {status\\== "started"?};

    %% Sorties
    \node[ok] (succ) at (7, 1.5) {[OK] Formater succès\\timestamp · source};
    \node[err] (fail) at (11, 3.5) {[ERR] Formater erreur\\details};

    %% Flèches
    \draw[arr] (sched.east) -- ++(0.8,0) |- (http.west);
    \draw[arr] (wh.east) -- ++(0.8,0) |- (http.west);
    \draw[arr] (http) -- (cond);
    \draw[arr] (cond) -- node[left, font=\scriptsize\bfseries] {OUI} (succ);
    \draw[arr] (cond) -- node[above, font=\scriptsize\bfseries] {NON} (fail);
  \end{tikzpicture}
  \caption{Workflow 04 — Déclencheur de collecte automatique}
  \label{fig:wf04}
\end{figure}

\subsection{Workflow 03 — Démo live (pipeline bout-en-bout)}

Ce workflow illustre le pipeline complet en 8 étapes en une seule exécution
manuelle.

\begin{figure}[H]
  \centering
  \begin{tikzpicture}[
    node distance=0.2cm,
    n/.style={draw=darkBlue!60, fill=pipColor, rounded corners=3pt,
              minimum width=2.6cm, minimum height=0.8cm,
              font=\scriptsize\bfseries, align=center, inner sep=4pt},
    arr/.style={-{Stealth[length=5pt]}, darkBlue!60, thick}
  ]
    \node[n] (n1) {$\triangleright$ Manual\\Trigger};
    \node[n, right=0.3cm of n1] (n2) {Set Params\\COBOL/FR};
    \node[n, right=0.3cm of n2] (n3) {POST /api\\ingest};
    \node[n, right=0.3cm of n3] (n4) {Wait 3s};
    \node[n, right=0.3cm of n4] (n5) {GET /api\\prospects};
    \node[n, below=0.6cm of n5] (n6) {Extract\\opp\_id};
    \node[n, left=0.3cm of n6] (n7) {POST /api/\\generate/msg};
    \node[n, left=0.3cm of n7] (n8) {Set:\\Résultat};

    \foreach \a/\b in {n1/n2, n2/n3, n3/n4, n4/n5, n5/n6, n6/n7, n7/n8}
      \draw[arr] (\a) -- (\b);

    %% Annotations
    \node[font=\scriptsize\itshape\color{espritGray}, align=center, below=0.15cm of n3]
      {BNP Paribas\\COBOL injecté};
    \node[font=\scriptsize\itshape\color{espritGray}, align=center, below=0.15cm of n7]
      {Groq LLM\\1--3s};
    \node[font=\scriptsize\itshape\color{darkGreen}, align=center, below=0.15cm of n8]
      {score, message,\\lead\_id};
  \end{tikzpicture}
  \caption{Workflow 03 — Démo live bout-en-bout (8 nœuds)}
  \label{fig:wf03}
\end{figure}

%% ═══════════════════════════════════════════════════════════════════════════════
\section{Conclusion}

Ce chapitre a présenté la conception détaillée de tous les composants~: les 4
collecteurs actifs, le pipeline d'ingestion en 6 étapes avec sa formule de scoring
pondérée, les deux classifieurs NLP indépendants, le générateur de messages LLM,
le schéma de base de données relationnel, et les 4 workflows n8n. Le chapitre
suivant présente la réalisation concrète et les résultats obtenus.
